% \section{Agent Prompts, Runner Scripts, and Logging Schema}
% \label{app:agent_scripts}
% \label{app:logging}

% We provide (i) baseline task prompts, (ii) the two-phase diagnosis/apply prompt templates, (iii) reference runner scripts (e.g., a Claude Code wrapper), and (iv) the logging schema used to record conversations, patches, deploy logs, and measurement artifacts.

% We provide (i) baseline task prompts, (ii) the two-phase diagnosis/apply prompt templates, (iii) reference runner scripts (e.g., a Claude Code wrapper), and (iv) the logging schema used to record conversations, patches, deploy logs, and measurement 
% artifacts.

% \section{Agent Prompts}
% \label{sec:prompts}
% \begin{tcblisting}{
%   title=Planning Agent Prompt,
%   colback=gray!10,
%   colframe=black,
%   coltitle=white,
%   colbacktitle=black,
%   fonttitle=\bfseries,
%   breakable,
%   enhanced,
%   boxed title style={
%     sharp corners,
%     boxrule=0pt
%   },
%   listing only
% }

% You are a Core Web Vitals optimization expert analyzing a $FRAMEWORK web application.

% ## Task: Optimize LCP, CLS, and INP for mobile and desktop

% ### Current Baseline Metrics:

% **Mobile:**
% - LCP (Largest Contentful Paint): ${LCP_MOBILE_MEDIAN}ms - Rating: ${LCP_MOBILE_RATING}
% - CLS (Cumulative Layout Shift): ${CLS_MOBILE_MEDIAN} - Rating: ${CLS_MOBILE_RATING}
% - INP (Interaction to Next Paint): ${INP_MOBILE_MEDIAN}ms - Rating: ${INP_MOBILE_RATING}

% **Desktop:**
% - LCP (Largest Contentful Paint): ${LCP_DESKTOP_MEDIAN}ms - Rating: ${LCP_DESKTOP_RATING}
% - CLS (Cumulative Layout Shift): ${CLS_DESKTOP_MEDIAN} - Rating: ${CLS_DESKTOP_RATING}
% - INP (Interaction to Next Paint): ${INP_DESKTOP_MEDIAN}ms - Rating: ${INP_DESKTOP_RATING}

% ### Your Task:  
% Analyze the codebase and create a detailed optimization plan to improve these Core Web Vitals metrics:
% - **LCP**: Optimize loading of main content elements
% - **CLS**: Prevent unexpected layout shifts during page load
% - **INP**: Improve responsiveness to user interactions

% ### Data Available:
% - repo/init_cwv.json: Contains detailed LCP element entries showing which elements are causing slow LCP
% - repo/: Complete source code for the application

% ### Output Instructions:
% - Read repo/init_cwv.json to understand which specific elements are causing LCP issues
% - Analyze the codebase to identify optimization opportunities
% - WRITE the plan to 'plan.md' in the current directory
% - Prioritize metrics with "Poor" or "Needs Improvement" ratings
% - Be specific about file paths and code changes
% - Use valid Markdown formatting
% - DO NOT modify any repository files (init_cwv.json or source code)
% - DO NOT create additional files or output to chat

% \end{tcblisting}


% \begin{tcblisting}{
%   title=Execution Agent Prompt,
%   colback=gray!10,
%   colframe=black,
%   coltitle=white,
%   colbacktitle=black,
%   fonttitle=\bfseries,
%   breakable,
%   enhanced,
%   boxed title style={
%     sharp corners,
%     boxrule=0pt
%   },
%   listing only
% }

% You are implementing Core Web Vitals optimizations for a ${FRAMEWORK} website.

% ### Current Baseline Metrics:

% **Mobile:**
% - LCP (Largest Contentful Paint): ${LCP_MOBILE_MEDIAN}ms - Rating: ${LCP_MOBILE_RATING}
% - CLS (Cumulative Layout Shift): ${CLS_MOBILE_MEDIAN} - Rating: ${CLS_MOBILE_RATING}
% - INP (Interaction to Next Paint): ${INP_MOBILE_MEDIAN}ms - Rating: ${INP_MOBILE_RATING}

% **Desktop:**
% - LCP (Largest Contentful Paint): ${LCP_DESKTOP_MEDIAN}ms - Rating: ${LCP_DESKTOP_RATING}
% - CLS (Cumulative Layout Shift): ${CLS_DESKTOP_MEDIAN} - Rating: ${CLS_DESKTOP_RATING}
% - INP (Interaction to Next Paint): ${INP_DESKTOP_MEDIAN}ms - Rating: ${INP_DESKTOP_RATING}

% ### The Plan:
% ${PLAN_CONTENT}

% ### Your Task:
% Execute the code modifications specified in the plan above to optimize CWV metrics (LCP, CLS, INP) for both mobile and desktop.

% ### ${FRAMEWORK}-Specific Considerations:
% - Work within the existing ${FRAMEWORK} architecture and patterns
% - Preserve all existing functionality and visible content

% ### Implementation Constraints:
% - Follow the plan and implement the changes
% - You can read files to get better understanding of the codebase
% - Do NOT edit init_cwv.json or configuration files
% - Do NOT remove pages or alter visible content/layout
% - Apply optimizations that work for both mobile and desktop viewports

% Focus on executing the concrete file modifications from the plan. Skip any analysis or documentation steps.


% \end{tcblisting}







% % \section{Internal Page Discovery and Deduplication Details}
% % \label{app:page_dedup}

% % We provide full details of sitemap parsing, bounded crawl, URL normalization rules, template canonicalization, and DOM-level deduplication, along with ablations on page caps and depth bounds.

% % ============================================================
% % Appendix: Internal Page Expansion (Sitemap/Crawl + Dedup)
% % ============================================================
\section{Framework Host Scripts}
\label{app:host-scripts}

We provide the complete set of deterministic host scripts used to deploy each framework locally during evaluation. Each script accepts a repository directory and log path, installs dependencies where necessary, and launches a local server on a fixed port. If no recognized entrypoint is found, the script exits with a non-zero status and the instance is excluded from evaluation.

\begin{lstlisting}[
  caption={\texttt{host\_static\_html.sh} -- Static HTML deployment.},
  label={lst:host-static},
  language=bash,
  basicstyle=\ttfamily\scriptsize,
  breaklines=true,
  numbers=left,
  numberstyle=\tiny\color{gray},
  xleftmargin=2em,
  framexleftmargin=1.5em
]
#!/usr/bin/env bash
set -euo pipefail
PORT="${PORT:-4000}"
REPO_DIR="${1:-.}"
LOG="${2:-/tmp/host_static_html.log}"
cd "$REPO_DIR"

if [[ -f index.html ]]; then
  echo "[static] Serving static HTML on port $PORT" \
    | tee -a "$LOG"
  exec python3 -m http.server "$PORT" \
    >>"$LOG" 2>&1
fi

echo "[static] ERROR: No index.html found" \
  | tee -a "$LOG"
exit 2
\end{lstlisting}

\begin{lstlisting}[
  caption={\texttt{host\_jekyll.sh} -- Jekyll deployment.},
  label={lst:host-jekyll},
  language=bash,
  basicstyle=\ttfamily\scriptsize,
  breaklines=true,
  numbers=left,
  numberstyle=\tiny\color{gray},
  xleftmargin=2em,
  framexleftmargin=1.5em
]
#!/usr/bin/env bash
set -euo pipefail
PORT="${PORT:-4000}"
REPO_DIR="${1:-.}"
LOG="${2:-/tmp/host_jekyll.log}"
cd "$REPO_DIR"

if [[ -f Gemfile ]]; then
  bundle install
  exec bundle exec jekyll serve \
    --host 0.0.0.0 --port "$PORT" >>"$LOG" 2>&1
fi

if [[ -f _config.yml ]]; then
  exec jekyll serve \
    --host 0.0.0.0 --port "$PORT" >>"$LOG" 2>&1
fi

if [[ -f index.html ]]; then
  exec python3 -m http.server "$PORT" >>"$LOG" 2>&1
fi

echo "[jekyll] ERROR: No recognized entrypoint" \
  | tee -a "$LOG"
exit 2
\end{lstlisting}

\begin{lstlisting}[
  caption={\texttt{host\_hugo.sh} -- Hugo deployment.},
  label={lst:host-hugo},
  language=bash,
  basicstyle=\ttfamily\scriptsize,
  breaklines=true,
  numbers=left,
  numberstyle=\tiny\color{gray},
  xleftmargin=2em,
  framexleftmargin=1.5em
]
#!/usr/bin/env bash
set -euo pipefail
PORT=${PORT:-8080}
REPO_DIR=${1:-.}
LOG=${2:-/tmp/host_hugo.log}
cd "$REPO_DIR"

for cfg in hugo.toml hugo.yaml hugo.yml \
           config.toml config.yaml config.yml; do
  if [ -f "$cfg" ]; then
    hugo server -p "$PORT" --bind 0.0.0.0 \
      >>"$LOG" 2>&1 &
    echo $! > /tmp/host_hugo.pid; exit 0
  fi
done

if [ -f public/index.html ]; then
  python3 -m http.server "$PORT" --directory public \
    >>"$LOG" 2>&1 &
  echo $! > /tmp/host_hugo.pid; exit 0
fi
if [ -f docs/index.html ]; then
  python3 -m http.server "$PORT" --directory docs \
    >>"$LOG" 2>&1 &
  echo $! > /tmp/host_hugo.pid; exit 0
fi
if [ -f index.html ]; then
  python3 -m http.server "$PORT" >>"$LOG" 2>&1 &
  echo $! > /tmp/host_hugo.pid; exit 0
fi

echo "[hugo] No recognized entrypoint" | tee -a "$LOG"
exit 2
\end{lstlisting}

\begin{lstlisting}[
  caption={\texttt{host\_hexo.sh} -- Hexo deployment.},
  label={lst:host-hexo},
  language=bash,
  basicstyle=\ttfamily\scriptsize,
  breaklines=true,
  numbers=left,
  numberstyle=\tiny\color{gray},
  xleftmargin=2em,
  framexleftmargin=1.5em
]
#!/usr/bin/env bash
set -euo pipefail
PORT=${PORT:-8080}
REPO_DIR=${1:-.}
LOG=${2:-/tmp/host_hexo.log}
cd "$REPO_DIR"

if [ -f package.json ]; then
  npm install --silent || true
  npx hexo server -p "$PORT" --silent \
    >>"$LOG" 2>&1 &
  echo $! > /tmp/host_hexo.pid; exit 0
fi

if [ -f index.html ]; then
  python3 -m http.server "$PORT" >>"$LOG" 2>&1 &
  echo $! > /tmp/host_hexo.pid; exit 0
fi

echo "[hexo] No recognized entrypoint" | tee -a "$LOG"
exit 2
\end{lstlisting}

\begin{lstlisting}[
  caption={\texttt{host\_react.sh} -- React deployment.},
  label={lst:host-react},
  language=bash,
  basicstyle=\ttfamily\scriptsize,
  breaklines=true,
  numbers=left,
  numberstyle=\tiny\color{gray},
  xleftmargin=2em,
  framexleftmargin=1.5em
]
#!/usr/bin/env bash
set -euo pipefail
PORT=${PORT:-8080}
REPO_DIR=${1:-.}
LOG=${2:-/tmp/host_react.log}
cd "$REPO_DIR"

if [ -f package.json ]; then
  npm install --silent || true
  export PORT="$PORT"
  npm start >>"$LOG" 2>&1 &
  echo $! > /tmp/host_react.pid; exit 0
fi

for vf in vite.config.ts vite.config.js; do
  if [ -f "$vf" ]; then
    npm install --silent || true
    npm run dev -- --host 0.0.0.0 --port "$PORT" \
      >>"$LOG" 2>&1 &
    echo $! > /tmp/host_react.pid; exit 0
  fi
done

for d in dist build; do
  if [ -f "$d/index.html" ]; then
    python3 -m http.server "$PORT" --directory "$d" \
      >>"$LOG" 2>&1 &
    echo $! > /tmp/host_react.pid; exit 0
  fi
done

if [ -f index.html ]; then
  python3 -m http.server "$PORT" >>"$LOG" 2>&1 &
  echo $! > /tmp/host_react.pid; exit 0
fi

echo "[react] No recognized entrypoint" | tee -a "$LOG"
exit 2
\end{lstlisting}

\begin{lstlisting}[
  caption={\texttt{host\_next.sh} -- Next.js deployment.},
  label={lst:host-next},
  language=bash,
  basicstyle=\ttfamily\scriptsize,
  breaklines=true,
  numbers=left,
  numberstyle=\tiny\color{gray},
  xleftmargin=2em,
  framexleftmargin=1.5em
]
#!/usr/bin/env bash
set -euo pipefail
PORT=${PORT:-8080}
REPO_DIR=${1:-.}
LOG=${2:-/tmp/host_next.log}
cd "$REPO_DIR"

if [ -f package.json ]; then
  npm install --silent || true
  npm run build --silent || true
  PORT="$PORT" npm run start >>"$LOG" 2>&1 &
  echo $! > /tmp/host_next.pid; exit 0
fi

for sub in website web; do
  if [ -f "$sub/package.json" ]; then
    (cd "$sub" && npm install --silent || true \
      && npm run build --silent || true \
      && PORT="$PORT" npm run start >>"$LOG" 2>&1 &)
    echo $! > /tmp/host_next.pid; exit 0
  fi
done

for d in out _next/static .; do
  if [ -f "$d/index.html" ] || [ -d "$d" ]; then
    python3 -m http.server "$PORT" --directory "$d" \
      >>"$LOG" 2>&1 &
    echo $! > /tmp/host_next.pid; exit 0
  fi
done

echo "[next] No recognized entrypoint" | tee -a "$LOG"
exit 2
\end{lstlisting}

\begin{lstlisting}[
  caption={\texttt{host\_vue.sh} -- Vue deployment.},
  label={lst:host-vue},
  language=bash,
  basicstyle=\ttfamily\scriptsize,
  breaklines=true,
  numbers=left,
  numberstyle=\tiny\color{gray},
  xleftmargin=2em,
  framexleftmargin=1.5em
]
#!/usr/bin/env bash
set -euo pipefail
PORT=${PORT:-8080}
REPO_DIR=${1:-.}
LOG=${2:-/tmp/host_vue.log}
cd "$REPO_DIR"

if [ -f package.json ]; then
  npm install --silent || true
  export PORT="$PORT"
  npm run serve >>"$LOG" 2>&1 &
  echo $! > /tmp/host_vue.pid; exit 0
fi

for vf in vite.config.ts vite.config.js; do
  if [ -f "$vf" ]; then
    npm install --silent || true
    npm run dev -- --host 0.0.0.0 --port "$PORT" \
      >>"$LOG" 2>&1 &
    echo $! > /tmp/host_vue.pid; exit 0
  fi
done

if [ -f dist/index.html ]; then
  python3 -m http.server "$PORT" --directory dist \
    >>"$LOG" 2>&1 &
  echo $! > /tmp/host_vue.pid; exit 0
fi

if [ -f index.html ]; then
  python3 -m http.server "$PORT" >>"$LOG" 2>&1 &
  echo $! > /tmp/host_vue.pid; exit 0
fi

echo "[vue] No recognized entrypoint" | tee -a "$LOG"
exit 2
\end{lstlisting}

\begin{lstlisting}[
  caption={\texttt{host\_express.sh} -- Express deployment.},
  label={lst:host-express},
  language=bash,
  basicstyle=\ttfamily\scriptsize,
  breaklines=true,
  numbers=left,
  numberstyle=\tiny\color{gray},
  xleftmargin=2em,
  framexleftmargin=1.5em
]
#!/usr/bin/env bash
set -euo pipefail
PORT=${PORT:-8080}
REPO_DIR=${1:-.}
LOG=${2:-/tmp/host_express.log}
cd "$REPO_DIR"

if [ -f package.json ]; then
  npm install --silent || true
  export PORT="$PORT"
  if npm run start --silent >/dev/null 2>&1; then
    npm run start >>"$LOG" 2>&1 &
    echo $! > /tmp/host_express.pid; exit 0
  fi
fi

for f in server.js app.js index.js; do
  if [ -f "$f" ]; then
    npm install --silent || true
    export PORT="$PORT"
    node "$f" >>"$LOG" 2>&1 &
    echo $! > /tmp/host_express.pid; exit 0
  fi
done

if [ -f backend/package.json ]; then
  (cd backend && npm install --silent || true \
    && export PORT="$PORT" \
    && npm run start >>"$LOG" 2>&1 &)
  echo $! > /tmp/host_express.pid; exit 0
fi

echo "[express] No recognized entrypoint" \
  | tee -a "$LOG"
exit 2
\end{lstlisting}

\begin{lstlisting}[
  caption={\texttt{host\_pelican.sh} -- Pelican deployment.},
  label={lst:host-pelican},
  language=bash,
  basicstyle=\ttfamily\scriptsize,
  breaklines=true,
  numbers=left,
  numberstyle=\tiny\color{gray},
  xleftmargin=2em,
  framexleftmargin=1.5em
]
#!/usr/bin/env bash
set -euo pipefail
PORT=${PORT:-8080}
REPO_DIR=${1:-.}
LOG=${2:-/tmp/host_pelican.log}
cd "$REPO_DIR"

if [ -f pelicanconf.py ]; then
  pip install -r requirements.txt || true
  pelican content || true
  python3 -m http.server "$PORT" --directory output \
    >>"$LOG" 2>&1 &
  echo $! > /tmp/host_pelican.pid; exit 0
fi

if [ -f publishconf.py ]; then
  pip install -r requirements.txt || true
  pelican content -s publishconf.py || true
  python3 -m http.server "$PORT" --directory output \
    >>"$LOG" 2>&1 &
  echo $! > /tmp/host_pelican.pid; exit 0
fi

if [ -f output/index.html ]; then
  python3 -m http.server "$PORT" --directory output \
    >>"$LOG" 2>&1 &
  echo $! > /tmp/host_pelican.pid; exit 0
fi

if [ -f index.html ]; then
  python3 -m http.server "$PORT" >>"$LOG" 2>&1 &
  echo $! > /tmp/host_pelican.pid; exit 0
fi

echo "[pelican] No recognized entrypoint" \
  | tee -a "$LOG"
exit 2
\end{lstlisting}

\begin{lstlisting}[
  caption={\texttt{host\_quarto.sh} -- Quarto deployment.},
  label={lst:host-quarto},
  language=bash,
  basicstyle=\ttfamily\scriptsize,
  breaklines=true,
  numbers=left,
  numberstyle=\tiny\color{gray},
  xleftmargin=2em,
  framexleftmargin=1.5em
]
#!/usr/bin/env bash
set -euo pipefail
PORT=${PORT:-8080}
REPO_DIR=${1:-.}
LOG=${2:-/tmp/host_quarto.log}
cd "$REPO_DIR"

if [ -f _quarto.yml ] || [ -f _quarto.yaml ]; then
  quarto render || true
  python3 -m http.server "$PORT" --directory _site \
    >>"$LOG" 2>&1 &
  echo $! > /tmp/host_quarto.pid; exit 0
fi

for d in _site docs .; do
  if [ -f "$d/index.html" ]; then
    python3 -m http.server "$PORT" --directory "$d" \
      >>"$LOG" 2>&1 &
    echo $! > /tmp/host_quarto.pid; exit 0
  fi
done

echo "[quarto] No recognized entrypoint" \
  | tee -a "$LOG"
exit 2
\end{lstlisting}

\begin{lstlisting}[
  caption={\texttt{host\_flask.sh} -- Flask deployment.},
  label={lst:host-flask},
  language=bash,
  basicstyle=\ttfamily\scriptsize,
  breaklines=true,
  numbers=left,
  numberstyle=\tiny\color{gray},
  xleftmargin=2em,
  framexleftmargin=1.5em
]
#!/usr/bin/env bash
set -euo pipefail
PORT=${PORT:-8080}
REPO_DIR=${1:-.}
LOG=${2:-/tmp/host_flask.log}
cd "$REPO_DIR"

if [ -f app.py ]; then
  pip install -r requirements.txt || true
  export FLASK_APP=app.py
  export FLASK_ENV=development
  export FLASK_RUN_PORT="$PORT"
  flask run --host=0.0.0.0 >>"$LOG" 2>&1 &
  echo $! > /tmp/host_flask.pid; exit 0
fi

if [ -f wsgi.py ]; then
  pip install -r requirements.txt || true
  export FLASK_APP=wsgi.py
  export FLASK_ENV=development
  export FLASK_RUN_PORT="$PORT"
  flask run --host=0.0.0.0 >>"$LOG" 2>&1 &
  echo $! > /tmp/host_flask.pid; exit 0
fi

if [ -f static/index.html ]; then
  python3 -m http.server "$PORT" --directory static \
    >>"$LOG" 2>&1 &
  echo $! > /tmp/host_flask.pid; exit 0
fi

echo "[flask] No recognized entrypoint" \
  | tee -a "$LOG"
exit 2
\end{lstlisting}
